\section{置換との併用}\label{節:置換と部分積分}

積分法は,問題ごとに1つだけ選ぶものではない.前章でも,根号を変数に換えた後で割り算を使った.では,置換と部分積分を組み合わせるとどうなるだろうか.ここでは
\[
 \int f(\log x)\dd x
\]
の形の積分を考える.対数の中身を1つの変数にする置換をすれば積が現れ,それを部分積分で計算できる.

まず,\,$f$ の中にある $\log x$ を1つの変数 $t$ にしたい.しかし,もとの積分には $\log x$ の微分 $\frac1x$ がないので,\,$f(t)\dd t$ とそのまま置き換えることはできない.そこで $t=\log x$ を $x=\e^t$ と解き直してやる.\,$\log x$ は狭義単調増加なので,逆向きの対応 $x=\e^t$ も1つに決まる.これを微分すると
\[
 \dd x=\e^t\dd t
\]
である.この微分の係数も含めて換えると,
\[
 \int f(\log x)\dd x=\int f(t)\e^t\dd t
\]
となる.置換によって $\dd x$ から $\e^t$ という因子が現れ,右辺は $f(t)$ と $\e^t$ の積になった.\,$\e^t$ は積分しても $\e^t$ のままなので,積分する側にうってつけである.\,$u=f(t)$,\,$v'=\e^t$ として部分積分すれば,
\[
 \int f(t)\e^t\dd t
 =f(t)\e^t-\int f'(t)\e^t\dd t
\]
となる.\,$f$ が整式や三角関数なら,前節までの方法で右辺の積分を求められる.最後に $t=\log x$ を代入して,もとの変数へ戻せばよい.

\begin{bookpoint}
\textbf{POINT}\par
$t=\log x$ と置いて $\dd x=\e^t\dd t$ とし,現れた積 $f(t)\e^t$ を部分積分してから $t=\log x$ に戻す.
\end{bookpoint}

\begin{example}{対数と三角関数}
$x>0$ において,\,$\int\sin(\log x)\dd x$ を求めよ.
\end{example}

\paragraphbf{方針}
$t=\log x$ と置くと $\int\e^t\sin t\dd t$ になる.これは前節の例題の形である.

\paragraphbf{解答}
$\frac x2\{\sin(\log x)-\cos(\log x)\}+C$

\paragraphbf{解法}
$t=\log x$ と置くと,\,$x>0$ だから $x=\e^t$ であり,\,$\dd x=\e^t\dd t$ となる.したがって
\[
 \int\sin(\log x)\dd x
 =\int\e^t\sin t\dd t
\]
である.前節の例題の結果を使うと,
\[
\begin{aligned}
 \int\sin(\log x)\dd x
 &=\frac12\e^t(\sin t-\cos t)+C\\
 &=\frac x2\{\sin(\log x)-\cos(\log x)\}+C
\end{aligned}
\]
を得る.最後の行では $\e^t=x$,\,$t=\log x$ を代入した.

検算では,外にある $x$ と,内側の $\log x$ の両方を微分する.
\[
\begin{aligned}
 &\frac{\dd}{\dd x}
 \left[\frac x2\{\sin(\log x)-\cos(\log x)\}\right]\\
 &\qquad=\frac12\{\sin(\log x)-\cos(\log x)\}
   +\frac x2\cdot\frac1x
     \{\cos(\log x)+\sin(\log x)\}\\
 &\qquad=\sin(\log x)
\end{aligned}
\]

\paragraphbf{手法の役割}
$t=\log x$ と置くだけでは計算は終わらず,置換後の積を部分積分で求め,最後に元の変数へ戻した.置換は合成関数の中身を1つの変数にする役割を,部分積分は $\e^t$ との積を処理する役割を担っている.途中の式の形に合わせて,次に使う方法を選べばよいのである.
